% throttling-qps.tex — Figura obligatoria 2 (R-30): throttling contra throughput.
%
% Lo que hay que saber para leer el eje (medido en F6.2, declarado en SPEC §4):
% `cpu_throttled_seconds` es la SUMA por CPU del tiempo que el cgroup tuvo hilos
% listos y no pudo correr. Por eso puede superar la duración de la ventana: con 4
% CPUs y una cuota de 1, cada segundo de reloj acumula hasta 4 de throttling. El
% eje NO se acota a 300 s.
%
% Además (F6.3): parte del throttling de MongoDB en el perfil C ocurre en reposo,
% sin benchmark (~27 s por ventana de 300 s). La línea de referencia lo marca.
%
% Un archivo de datos por motor y concurrencia (report/datos/throttling_*.dat),
% para no depender de filtros de pgfplots.

\begin{figure}[htbp]
  \centering
  \begin{tikzpicture}
    \begin{axis}[
        xlabel={Throttling de CPU en la ventana de 300~s (s)},
        ylabel={Throughput (consultas/s)},
        ymode=log, log basis y=10,
        xmin=0,   % el throttling no puede ser negativo
        enlarge x limits={upper, value=0.05},
        legend pos=south west, legend cell align=left, legend columns=2,
        height=7.5cm, width=\linewidth,
      ]
      \addplot[only marks, mark=*,       mark size=1.6pt, blue!65!black]
        table [x=throttling_s, y=qps] {datos/throttling_mongodb_c1.dat};
      \addplot[only marks, mark=*,       mark size=2.6pt, blue!65!black]
        table [x=throttling_s, y=qps] {datos/throttling_mongodb_c10.dat};
      \addplot[only marks, mark=*,       mark size=3.6pt, blue!65!black]
        table [x=throttling_s, y=qps] {datos/throttling_mongodb_c50.dat};
      \addplot[only marks, mark=square*, mark size=1.6pt, red!65!black]
        table [x=throttling_s, y=qps] {datos/throttling_postgresql_c1.dat};
      \addplot[only marks, mark=square*, mark size=2.6pt, red!65!black]
        table [x=throttling_s, y=qps] {datos/throttling_postgresql_c10.dat};
      \addplot[only marks, mark=square*, mark size=3.6pt, red!65!black]
        table [x=throttling_s, y=qps] {datos/throttling_postgresql_c50.dat};
      \legend{MongoDB c=1, MongoDB c=10, MongoDB c=50,
              PostgreSQL c=1, PostgreSQL c=10, PostgreSQL c=50}
      % F6.3: línea base de throttling en reposo de MongoDB en el perfil C.
      \draw[dashed, gray!70] ({axis cs:26.7,0}|-{rel axis cs:0,0})
                          -- ({axis cs:26.7,0}|-{rel axis cs:0,1});
      \node[gray!70, font=\tiny, anchor=south west, rotate=90]
        at ({axis cs:26.7,0}|-{rel axis cs:0,0.05}) {MongoDB · C en reposo (F6.3)};
    \end{axis}
  \end{tikzpicture}
  \caption[Throttling contra throughput]{%
    \textbf{Throttling de CPU contra throughput.} Cada punto es una celda (media de
    sus 3 repeticiones); el tamaño de la marca crece con la concurrencia.
    \textbf{El eje horizontal supera los 300~s de la ventana} porque
    \texttt{cpu\_throttled\_seconds} suma el tiempo de cada CPU en la que el cgroup
    tenía hilos listos (medido en F6.2, declarado en SPEC §4); el máximo observado
    es de \num{2086}~s. La línea vertical marca el throttling que MongoDB acumula
    \textbf{en reposo} en el perfil C, sin benchmark (F6.3): a su izquierda no hay
    contención atribuible a la carga. Datos:
    \texttt{report/datos/throttling\_*.dat}.}
  \label{fig:throttling-qps}
\end{figure}
